\documentclass{situationist-preprint}

\title{Proudhon's Paradox}
\msid{POL-06-b}
\msfamily{POL-06}
\msdiscipline{Political Philosophy}
\mstype{Consequence}
\mskeywords{Proudhon; property; theft; possession; Marx; self-undermining theses}

\begin{document}
\maketitle

\begin{abstract}
Proudhon held that property is theft, and later defended property as a guarantee of liberty. The pairing is often presented as a paradox. We derive its resolution from the conclusion of POL-06: since theft presupposes property, the thesis that all property is theft is self-undermining, and every true thesis in its neighborhood condemns some property in the name of other property rights. It follows that Proudhon's two claims are consistent only if they concern different holdings, and that the interesting content of the first is not that property is theft but a claim about which holdings are rightful. That is a question in the theory of justice, not a refutation of property.
\end{abstract}

\section{The paradox}

In 1840 Proudhon answered the question of his title, what property is, with the thesis that it is theft \citep{Proudhon1840}. In his posthumous \emph{Th\'eorie de la propri\'et\'e} he defended property as a counterweight to the power of the state and a condition of individual liberty \citep{Proudhon1866}. Read together, the claims seem to say that one and the same institution is a crime and a safeguard of freedom. Marx, in his obituary letter on Proudhon, pressed the first claim from the other side: theft, as a forcible violation of property, presupposes property, so that the slogan loses itself in confused notions of what bourgeois property is \citep{Marx1865}.

POL-06-a proves Marx's point. Given that whatever is stolen is something to which someone else has a right (Assumption~2.3 of POL-06-a), and that rightful ownership is a form of holding (Assumption~2.4), the thesis that every holding is a theft is false as soon as anyone holds anything (Theorem~3.1 of POL-06-a), and can be true only in a world in which no theft occurs (Corollary~3.3).

\section{The resolution}

\begin{proposition}\label{p:split}
Any consistent pair of theses, one condemning property as theft and the other defending property, must concern different classes of holdings, and the first must presuppose that some holdings in the second class are rightful.
\end{proposition}

\begin{proof}
If both theses concerned the same holdings, one would say that they are all thefts and the other that some of them are rightful; by Definition~2.2 of POL-06-a, a theft is a holding that is not rightful, so the pair would be inconsistent. So the first thesis concerns a restricted class. If any holding in that class is a theft, then by Corollary~3.2 of POL-06-a some holding is not a theft, and by Assumption~2.3 of POL-06-a the agent wronged by the theft has a right to the thing, which by Assumption~2.4 is a rightful holding outside the condemned class.
\end{proof}

Proudhon himself supplies the division. What he condemned in 1840 was property as the right to an income without labor, rent, interest and profit, which he called the \emph{droit d'aubaine}; what he did not condemn was \emph{possession}, the holding of what one occupies and works \citep{Proudhon1840}. His thesis, so read, is that holdings of the first kind take from those who have a right to the product of their labor. By Proposition~\ref{p:split}, that thesis condemns some property in the name of other property rights, the rights of workers to what they produce. It is no longer a thesis about property as such.

\begin{corollary}\label{c:content}
The content of the restricted thesis is a claim about which holdings are rightful, and its truth cannot be settled by the concept of property.
\end{corollary}

\begin{proof}
By Proposition~\ref{p:split}, the restricted thesis distinguishes rightful from wrongful holdings. Whether income from ownership is wrongful depends on what persons are entitled to, which is not fixed by the definition of a holding or a right (Definition~2.1 of POL-06-a).
\end{proof}

Corollary~\ref{c:content} places Proudhon's question where it belongs: in the theory of just holdings. Whether income from capital is rightful is a question for the principles of POL-01 and for the assessment of markets in POL-12, and the answer of POL-12, that markets are just when their outcomes satisfy those principles, is neither Proudhon's nor his opponents'. The later defense of property is then not a reversal but the other half of the same view: possession is a guarantee of liberty, and certain property titles are not. Both halves are rights-based judgements of the kind whose general structure \citet{Waldron1988} examines: they differ in which rights they take holdings to answer to, not in whether holdings answer to rights.

\section{Objection}

\begin{objection}
The slogan was rhetoric, and to read it as a thesis is to miss its point.
\end{objection}

\begin{reply}
The point of the slogan was to borrow the condemnation that attaches to theft. The borrowing is legitimate only if the condemned holdings wrong someone with a right, and that is what Proposition~\ref{p:split} makes explicit. The rhetoric works only in virtue of the presupposition it seems to deny.
\end{reply}

\section{Conclusion}

Proudhon's paradox dissolves once the conclusion of POL-06 is applied to it. The unrestricted thesis is self-undermining; the restricted thesis is consistent with the defense of property because it is a claim about which property is rightful. What remains of the slogan is a question about just holdings, to be answered by a theory of justice.

\begin{thebibliography}{9}
\bibitem[Marx(1865)]{Marx1865} Marx, K. (1865). On Proudhon [letter to J.~B. Schweitzer, 24 January 1865]. In K. Marx and F. Engels, \emph{Collected Works}, vol.~20. Lawrence \& Wishart, 1985.
\bibitem[Proudhon(1840)]{Proudhon1840} Proudhon, P.-J. (1840). \emph{Qu'est-ce que la propri\'et\'e?} Trans. as \emph{What Is Property?}, ed. D.~R. Kelley and B.~G. Smith, Cambridge University Press, 1994.
\bibitem[Proudhon(1866)]{Proudhon1866} Proudhon, P.-J. (1866). \emph{Th\'eorie de la propri\'et\'e}. Lacroix, Verboeckhoven.
\bibitem[Waldron(1988)]{Waldron1988} Waldron, J. (1988). \emph{The Right to Private Property}. Clarendon Press.
\end{thebibliography}

\end{document}
